% ════════════
% RDI_Working_Paper.tex
% Relational Dignity Infrastructure: The Human Layer Required to Make Material Housing Inhabitable
% Material Dignity Infrastructure Working Paper 4
% ════════════

\documentclass[12pt, letterpaper]{article}

% ── PACKAGES ──
\usepackage[left=0.7in, right=0.7in, top=1in, bottom=1in, headheight=14pt]{geometry}
\usepackage[expansion=false]{microtype}
\usepackage{textcomp}
\usepackage{setspace}
\usepackage{booktabs}
\usepackage{tabularx}
\usepackage[titles]{tocloft}
\usepackage{tabularray}
\usepackage{longtable}
\usepackage{array}
\usepackage{multirow}
\usepackage{hyperref}
% natbib not used — manual thebibliography
\usepackage{amsmath}
\usepackage{amssymb}
\usepackage{parskip}
\usepackage{titlesec}
\usepackage{fancyhdr}
\usepackage[table]{xcolor}
\usepackage{caption}
\usepackage{footnote}
\usepackage{enumitem}
\usepackage{tikz}
\usepackage{needspace}
\usepackage{etoolbox}
\usepackage{float}

% ── COLOR SYSTEM ───
\definecolor{auditblue}{HTML}{003366}
\definecolor{rowgray}{HTML}{F5F5F5}

% ── GLOBAL SPACING ───
\setstretch{1.4}
\setlength{\parindent}{0pt}
\setlength{\parskip}{8pt}
\hyphenpenalty=1000

% ── SECTION FORMATTING. ZERO BOLDING IN TEXT. ───
\titleformat{\section}
  {\normalfont\large\color{auditblue}}{\thesection.}{0.5em}{}
\titleformat{\subsection}
  {\normalfont\normalsize\color{auditblue}}{\thesubsection.}{0.5em}{}
\titlespacing*{\section}{0pt}{18pt}{6pt}
\titlespacing*{\subsection}{0pt}{12pt}{4pt}

% ── HEADER / FOOTER ───
\pagestyle{fancy}
\fancyhf{}
\rhead{\small\textit{Relational Dignity Infrastructure (RDI Working Paper 4)}}
\lhead{\small\textit{Working Paper, May 2026}}
\cfoot{\thepage}
\renewcommand{\headrulewidth}{0.4pt}

% ── TABLE SETTINGS ───
\renewcommand{\arraystretch}{1.3}

% ── ORPHAN / WIDOW CONTROL ───
\clubpenalty=10000
\widowpenalty=10000
\displaywidowpenalty=10000

% ── BIBLIOGRAPHY ORPHAN PREVENTION ───
\AtBeginEnvironment{thebibliography}{\interlinepenalty=10000}

% ── HYPERREF CONFIGURATION ───
\hypersetup{
  colorlinks=true,
  linkcolor=black,
  citecolor=black,
  urlcolor=black,
  pdftitle={Relational Dignity Infrastructure: The Human Layer Required to Make Material Housing Inhabitable},
  pdfauthor={DiBella, C.J.},
  pdfsubject={Housing First. Ontological Security. Identity Capital. Homelessness Recovery.},
  pdfkeywords={Relational Dignity Infrastructure, ontological security, identity capital, Housing First retention, Dunbar scale, Pod Steward, peer support, social network preservation, homelessness recovery, material dignity infrastructure, warm offer, institutional objectification, street-to-home pipeline}
}
\setcounter{tocdepth}{2}

\begin{document}
\captionsetup[table]{labelfont=normalfont, labelsep=period, skip=6pt}

% ── FRONT MATTER ───
\begin{titlepage}
\begin{center}
\setstretch{1.6}
{\LARGE \textbf{Relational Dignity Infrastructure}}\\[6pt]
{\large The Human Layer Required to Make Material Housing Inhabitable}\\[1.0cm]
{\normalsize \textbf{Charles J. DiBella}}\\[0.01cm]
{\normalsize Principal Systems Architect}\\[0.01cm]
{\normalsize Working Paper --- May 2026}\\[0.5cm]
\textbf{ABSTRACT}
\vspace{0.5cm}
\end{center}
\begin{minipage}{0.95\textwidth}
\setstretch{1.3}
\noindent

Material housing provision does not produce human recovery. Finland's Housing First infrastructure achieves eighty- to ninety-percent two-year retention; the United States achieves below forty percent. The difference is not the housing unit. It is the ambient relational environment within and surrounding that unit at the moment of occupancy. This paper defines Relational Dignity Infrastructure (RDI) as the systematically designed, continuously operating social environment that produces ontological security, identity capital accumulation, and the functional capacity for reintegration that material shelter alone cannot generate. RDI is positioned as the third operational layer of a three-part dignity stack: Material Dignity Infrastructure delivers biological survival conditions; RDI delivers the relational recognition and social continuity required for selfhood persistence; identity capital accumulation delivers the agentive selfhood required for reintegration into civic life. Drawing on Giddens's ontological security framework, Padgett's social-material conjunction findings, Cote and Levine's identity capital theory, and the 2025 PLOS ONE longitudinal evidence on Housing First outcomes, the paper establishes RDI as a formal theoretical construct and maps its five production conditions onto the existing MDI tower architecture: Dunbar-scale pod environments, the live-in Pod Steward role, peer support workforce deployment, the unconditional Ground Floor Commons, and the Ontological Permanence process specified in MDI-C. Two integration gaps requiring specification are identified: the relational intake protocol and the social network transition protocol. Five RDI verification metrics are proposed as leading indicators of the eight binary material metrics that constitute the MDI Singular Prototype Threshold.

\end{minipage}
\end{titlepage}

\pagenumbering{arabic}

% ── DOCUMENT METADATA ────
\newpage
\section*{Document Metadata}

{\footnotesize\setstretch{1.25}

\noindent\textbf{Keywords}\\[1pt]
Relational Dignity Infrastructure, ontological security, identity capital, Housing First retention, Dunbar scale, Pod Steward, peer support, social network preservation, homelessness recovery, material dignity infrastructure, warm offer, institutional objectification, street-to-home pipeline

\vspace{8pt}
\noindent\textbf{JEL Classification}
\begin{itemize}[nosep, before={\vspace{-8pt}}, leftmargin=2em]
    \item I14. Health and Inequality
    \item I38. Government Policy, Provision and Effects of Welfare Programs
    \item H75. State and Local Government, Health, Education, Welfare
    \item Z13. Economic Sociology, Economic Anthropology, Social and Economic Stratification
    \item R21. Urban, Rural, Regional, Real Estate, and Transportation Economics: Housing Demand
    \item J15. Economics of Minorities, Races, Indigenous Peoples, and Immigrants; Non-labor Discrimination
\end{itemize}

\vspace{8pt}
\noindent\textbf{Target eJournals}
\begin{itemize}[nosep, before={\vspace{-8pt}}, leftmargin=2em]
    \item Housing and Residential Economics eJournal
    \item Public Health Law and Policy eJournal
    \item State and Local Government eJournal
    \item Economic Sociology eJournal
    \item Social Capital, Networks \& Trust eJournal
    \item Poverty, Income Distribution \& Social Protection eJournal
\end{itemize}

\vspace{8pt}
\noindent\textbf{Statement of Necessity}\\[1pt]
The homelessness policy literature has produced two decades of Housing First evidence demonstrating that permanent housing with wraparound services outperforms treatment-first models. What the literature has not produced is a formal theoretical account of what the relational environment within Housing First infrastructure must contain in order to generate the outcomes that Housing First promises. The result is a gap between architectural intention and human outcome that institutional service systems fill, by default, with transactional case management, behavioral compliance monitoring, and the objectification of residents as service recipients rather than as persons under reconstruction. This paper argues that this gap is not a service delivery problem. It is a theoretical problem: the field lacks a named, formally defined concept for the relational layer that material housing must contain. Without a concept, the layer cannot be designed. Without design, it cannot be replicated. Without replication, Housing First produces the forty-percent retention rates observed in the American system rather than the eighty- to ninety-percent rates observed in the Finnish system. Relational Dignity Infrastructure names the layer, defines its production conditions, and maps its implementation within the MDI tower architecture as a replicable, measurable infrastructure component.

\vspace{8pt}
\noindent\textbf{Author's Note}\\[1pt]
This paper is the fourth working paper in the Material Dignity Infrastructure series. MDI-A (SSRN 5968756) established the distributed stewardship governance architecture. MDI-B (SSRN 6211658) diagnosed structural misalignment and surplus shelter capacity at the national scale. MDI-C (SSRN 6579600) delivered the full industrial specification of the Los Angeles Metropolitan Stabilization pipeline. This paper establishes the relational architecture that must operate within that physical environment for the pipeline to produce durable human outcomes rather than measured occupancy.

\vspace{8pt}
\noindent\textbf{Suggested Citation}\\[1pt]
DiBella, C.J. (2026). Relational Dignity Infrastructure: The Human Layer Required to Make Material Housing Inhabitable. Material Dignity Infrastructure Working Paper 4. Working Paper, May 2026.

}

% ── TABLE OF CONTENTS ────
\newpage
\begingroup
\setstretch{1.05}
\setlength{\cftbeforesecskip}{2pt}
\setlength{\cftbeforesubsecskip}{-2pt}
\tableofcontents
\endgroup

% ── BODY ────
\newpage
\setstretch{1.35}

% ═════════════
\clearpage
\section{The Retention Problem and Its Missing Layer}

Finland reduced chronic homelessness by sixty-eight percent. Its Housing First program sustains eighty- to ninety-percent two-year housing retention. The United States operates Housing First programs in every major metropolitan area and achieves below forty percent retention. The housing unit is not the variable that differentiates these outcomes. The relational environment within and surrounding that unit is.

American Housing First implementation places individuals directly from street-level acute physiological crisis into permanent apartments. The MDI series has established the biological case for why this sequencing fails: sleep-destroyed, cortisol-flooded, prefrontal-cortex-compromised individuals cannot maintain housing tenure because the neurological equipment required to manage a lease has been systematically destroyed by years of street exposure. The biological argument is necessary but insufficient. Even individuals who complete metabolic stabilization and arrive in a permanent unit with intact neurological function fail at rates that purely biological models cannot explain.

The evidence from Padgett's longitudinal study of Housing First participants, from the 2025 PLOS ONE psychiatric outcome data, and from two decades of Housing First implementation science identifies a consistent predictor of two-year retention that operates independently of neurological status: the quality of the relational environment within the housing setting. Individuals embedded in stable relational networks, recognized as specific persons by the people who share their residential environment, and holding roles within their residential community that generate identity continuity retain housing at rates approaching the Finnish benchmark. Individuals placed in housing without these conditions retain at American rates regardless of the quality of the unit.

The relational environment is not a wraparound service. It is not case management, peer support programming, or community building activity. It is an ambient, structural property of the housing environment itself: the continuous background condition against which individual recovery either occurs or fails to occur. This paper names that property Relational Dignity Infrastructure and establishes its formal theoretical standing, its five production conditions, and its implementation within the MDI tower architecture.

% ═════════════
\clearpage
\section{Theoretical Foundations}

\subsection{Ontological Security: The First Requirement}

Anthony Giddens defined ontological security as the confidence most human beings have in the continuity of their self-identity and in the constancy of the surrounding social and material environments. It is not psychological stability in the clinical sense. It is the existential foundation beneath daily action: the taken-for-granted background of recognizable routine, stable social relationships, and spatial predictability that allows a person to plan, hope, and act in time without constant vigilance against environmental disruption.

Street homelessness systematically destroys all four of the conditions Giddens identifies as constitutive of ontological security. Routine collapses under the pressure of daily survival: food, water, shelter, and safety absorb all cognitive bandwidth, leaving nothing for the forward planning that routine requires. Social continuity fractures because the encampment environment, despite its real community bonds, is perpetually disrupted by sweeps, deaths, relocations, and institutional interventions. Spatial control disappears when any surface a person occupies is claimable by police at any moment. Identity continuity, which Giddens identifies as the deepest requirement, dissolves when every institutional contact confirms the person as a category of social problem rather than as a specific named individual with a particular history.

Padgett's 2007 qualitative follow-up of Tsemberis's Housing First study documented the social-material conjunction directly: stable housing was necessary but not sufficient for ontological security recovery. The locked door, the private bath, and the acoustic isolation together formed the material platform that created the conditions under which ontological security became possible. What converted that possibility into reality was the presence of social relationships within the housing setting that provided ongoing recognition, continuity, and witness. Housing First succeeds where it succeeds because it provides both layers. It fails where it fails because material provision without relational infrastructure is not Housing First. It is warehousing in a better building.

\subsection{The PLOS ONE Evidence (2025)}

The 2025 PLOS ONE longitudinal study of Housing First outcomes across five national systems identified two independent predictors of two-year retention that held across all five systems after controlling for diagnosis, substance use status, service intensity, and unit quality. The first was ontological security score at 90 days post-housing, which aligns with the biological decompression logic established in MDI-C. The second was social network density at 90 days: the number of named, reciprocal relationships the resident maintained within their housing environment. Residents with three or more named in-building peer contacts at 90 days retained housing at rates twenty-three percentage points higher than residents with fewer than three contacts, controlling for all other variables. The relationship was dose-dependent: each additional named peer contact added approximately four percentage points to two-year retention probability.

This finding establishes social network density as a measurable, actionable infrastructure target. It is not a soft outcome or an aspirational condition. It is a leading indicator of retention, the lagging material metric that the MDI Singular Prototype Threshold measures.

\subsection{Identity Capital: The Accumulation Mechanism}

Cote and Levine's identity capital theory positions selfhood not as a fixed property but as an accumulation: a portfolio of commitments, roles, relationships, and competencies that the individual builds over time through engagement with social environments. Chronic street homelessness does not merely fail to produce identity capital; it actively destroys it. The chronically unhoused person loses employment history, social roles, community membership, educational credentials, and the narrative continuity through which identity capital is maintained and communicated. Klodnick and Samuels's 2020 longitudinal work on identity development following housing entry documents that the period immediately post-housing is the critical window for identity capital reconstruction: residents who engage in identity-generating activities, including employment, education, civic participation, and community governance, within the first year of housing show dramatically better five-year outcomes than residents who do not, even controlling for all clinical variables.

The mechanism is accumulation, not restoration. Identity capital is not recovered; it is built from what remains after street exposure. The institutional environment within the housing setting determines whether the four raw materials for accumulation are structurally available: roles, recognition, relationships, and meaningful activity. An environment that provides only monitoring, services, and compliance requirements does not provide those raw materials. It provides a more comfortable version of the objectifying institutional relationship the person has experienced throughout their history with service systems.

\subsection{The Objectification Failure}

The sociological literature on total institutions (Goffman, 1961) and on institutional capture mechanisms within the contemporary homeless service system (Roebuck et al., 2022) identifies a consistent structural pattern: service systems that process large numbers of socially marginal individuals under financial incentives tied to volume rather than outcomes systematically develop organizational cultures that objectify service recipients. Objectification is not cruelty. It is the predictable institutional response to the conflict between operational efficiency and personalized relationship: when managing populations at scale, the individual becomes a case, the case becomes a category, and the category receives standardized treatment that cannot be modified by the specific reality of any individual person.

The relational damage this produces is not primarily emotional. It is structural. Objectification eliminates social recognition itself: the condition of being known as a specific person, which Giddens identifies as constitutive of ontological security and which the PLOS ONE evidence identifies as the predictor of retention. A person who is known only as a client, a case number, or a service recipient cannot build the identity capital that Klodnick and Samuels identify as predictive of five-year housing stability. The objectifying environment is, in this precise sense, a hostile environment for recovery regardless of the material quality of the housing it provides.

% ═════════════
\clearpage
\section{Relational Dignity Infrastructure: Formal Definition}

Relational Dignity Infrastructure (RDI) is the systematically designed, continuously operating social environment that generates ontological security, sustains identity capital accumulation, and maintains the social recognition of each resident as a specific named person with a particular history. These are conditions that material housing provision alone cannot produce and that are required for material housing to achieve its intended outcome of durable residential stability.

RDI is infrastructure in the technical sense. Like material infrastructure, it must be designed to specification, built into the operational architecture of the housing environment, continuously maintained, and measured against defined performance thresholds. It is not a program, a service modality, or a staff training philosophy. It is an ambient, structural property of the housing environment: the relational quality that is present or absent regardless of whether any specific interaction is occurring at any given moment.

RDI differs from existing concepts in the homelessness literature in three specific ways. It is structural, not programmatic: RDI is a property of the built and social environment, not an activity delivered at scheduled times. It is ambient, not episodic: RDI operates continuously as the background condition of daily residential life, not during formal service contacts. It is measurable, not aspirational: RDI produces specific, quantifiable outcomes, including named peer contact counts, pod council attendance, voluntary engagement rates, and steward recognition indices, each of which functions as a leading indicator of the lagging material retention metrics.

% ═════════════
\clearpage
\section{The Five RDI Production Conditions}

Five structural conditions are required to produce RDI. Each is already present in the MDI tower architecture, either explicitly specified or implied by the design logic. This paper makes their relational function explicit and positions them as the RDI layer within the MDI operational framework.

\subsection{Dunbar-Scale Social Environment}

Human beings cannot maintain stable recognition-based social relationships with more than approximately 150 individuals simultaneously. Beyond this cognitive ceiling, social environments produce anonymity as their default condition: individuals cease to be recognizable as specific persons and become background figures in an undifferentiated crowd. Anonymity is the social correlate of ontological insecurity; it is the environmental condition under which the continuous background of recognition that Giddens requires is absent.

The MDI tower's partition into 13 pods of approximately 154 residents each is the primary RDI structural condition. The Dunbar pod creates the scale within which every resident can be known by name, history, and situation by at least some significant portion of their daily social environment. The pod does not guarantee relational recognition; it creates the scale condition under which relational recognition becomes cognitively possible. What converts that possibility into reality are the remaining four production conditions.

\subsection{The Pod Steward: Continuous Witness}

The live-in Pod Steward is the RDI delivery mechanism that most directly addresses the objectification failure documented in the service system literature. The steward's defining functional characteristic is not their role or title. It is their residential continuity: they live in the pod. They are present not during scheduled hours but in the morning when residents move through the kitchen, at irregular hours when crises occur, across the full texture of daily life within which relational knowledge of specific persons accumulates. A steward who has lived in Pod 7 for six months knows who takes their coffee at 6 a.m., who has been withdrawing for three days, who had a difficult night last Tuesday, and who received good news yesterday. This ambient knowledge, qualitatively different from anything a case file can contain, is what distinguishes continuous relational presence from scheduled service delivery.

Three design features make the steward's RDI function possible. First, residential continuity: the live-in model ensures relational knowledge accumulates across the full texture of daily life rather than through formal contact windows. Second, non-clinical role definition: the steward's explicit exclusion from clinical intervention and security enforcement removes the authority asymmetries that make genuine horizontal relationship impossible in service system contexts. Third, compensation as relational signal: sovereign residency and dining access rather than payroll wage positions the steward as a community member who earned their place through service, not an institutional representative dispatched to manage residents.

\subsection{Peer Support Workforce: The Relational Bridge}

The MDI staffing model's 22.0 FTE Peer Support Specialist allocation is the largest single workforce category in the tower. Peer specialists hold the unique relational position of being simultaneously workers within the service system and members of the community being served. Their lived experience of homelessness, addiction, or psychiatric crisis produces a recognition of commonality that no amount of professional training can replicate in workers without that history. The peer specialist encounter is structurally different from the worker-client encounter because it does not depend on institutional authority for its legitimacy. It depends on shared experience.

The peer workforce's RDI function operates across three domains. In the pre-housing phase, peer specialists build the relational trust during outreach that makes the warm offer credible, because the offer works when it comes from a relationship rather than an institution. In the tower environment, peer specialists provide non-supervisory ambient presence: relational availability that signals a different category of social encounter than any authority-bearing role can provide. In de-escalation contexts, peer specialists address the relational crises, including dignity violations, identity threats, and disconnection, that precede and generate the behavioral incidents that clinical de-escalation protocols address after the fact.

\subsection{The Ground Floor Commons: Unconditional Relational Access}

The MDI Ground Floor Commons operates 24 hours a day, 365 days a year, with no intake requirement and no behavioral test for access. Its unconditional access model is an RDI principle made physical: the recognition of the person as a person, prior to any assessment of their needs, compliance history, or service eligibility, is the foundational relational act. The commons does not offer services on the condition that the person accepts a relational frame determined by institutional authority. It offers presence on the condition of nothing.

The commons also functions as a civic de-stigmatization environment. Its permeability to community members, neighbors, and non-residents alongside tower residents dissolves the institutional boundary that marks housed individuals as categorically different, classifying them as homeless people in a homeless facility. For residents whose sense of civic belonging has been destroyed by years of institutional stigmatization, this dissolution of categorical boundary is a relational repair mechanism. The commons is not the most intensive RDI delivery environment in the MDI architecture. It is the first one: the place where the relational quality of the MDI system is first encountered, before any decision about housing has been made.

\subsection{Recurring Community Practice: The Pod Council and Common Kitchen}

Shared history is the relational mechanism through which a group of people occupying the same space becomes a community rather than a collection of individuals who happen to be co-resident. Shared history requires recurring shared practice: activities that occur regularly, that all members participate in or could participate in, and that generate the accumulation of common reference points, mutual expectations, and collective memory through which community identity forms.

The MDI architecture specifies two recurring community practices at the pod level: the weekly Pod Council and the pod common kitchen. The Pod Council is the most important structured RDI mechanism in the tower because it is identity-generating as well as history-generating: it positions every resident as a citizen of the pod community with a voice that produces effects on the shared environment. The resident who attends, speaks, and sees their suggestion acted on is doing identity reconstruction work that no clinical intervention can substitute for. The common kitchen generates daily encounters through which casual, non-instrumental relational knowledge accumulates. Both mechanisms require no clinical staff to operate. They require only that the meeting space and kitchen be present as physical infrastructure and that the Pod Steward facilitate rather than direct the operation.


% ═════════════
\clearpage
\section{Integration with MDI Architecture: The Process 5 Mapping}

The MDI-C industrial specification, which formally establishes the MDI five-process parallel architecture, names the fifth process Ontological Permanence. Its component elements are the Dunbar pod structure, STC 65 acoustic isolation, the Pod Steward, the pod common kitchen, biophilic nodes, and the Digital Access Center. Process 5 holds that these elements together rebuild the neurological conditions for selfhood destroyed by street life, positioning it as the closest existing MDI language to what RDI theory formally names. The RDI framework is the theoretical architecture that Process 5 implies but does not fully develop.

The mapping between Process 5 elements and RDI production conditions is direct:

{\singlespacing
\begin{tabularx}{\textwidth}{@{}lX@{}}
\toprule
\rowcolor{rowgray}
Process 5 Element & RDI Production Condition \\
\midrule
Dunbar pod structure & Dunbar-scale social environment; recognition at human cognitive scale \\
\rowcolor{rowgray}
STC 65 acoustic sanctuary & Biological decompression prerequisite; sensory precondition for relational engagement \\
Pod Steward & Continuous witness; cross-time relational memory; non-conditional recognition \\
\rowcolor{rowgray}
Pod common kitchen & Recurring community practice; shared history generation; casual relational encounter \\
Biophilic nodes & Environmental downregulation; sensory context supporting presence \\
\rowcolor{rowgray}
Digital Access Center & Identity capital accumulation; narrative expansion; future-oriented engagement \\
\bottomrule
\end{tabularx}\par
\vspace{8pt}
\noindent\textit{Table 1. Process 5 Element to RDI Production Condition Mapping.}
\par}

This mapping establishes that the RDI layer is not an addition to the MDI architecture. It is a formal theoretical account of what the MDI architecture already contains. What RDI theory adds is the explicit naming and measurement of the relational function: a precision that makes design, staffing specification, training, and replication possible in ways that the current architectural language does not yet enable.

% ═════════════
\clearpage
\section{Integration Gaps: Two Specifications Required}

Reading the MDI architecture through the RDI framework identifies two integration gaps where the relational infrastructure is not yet fully specified.

\subsection{The Relational Intake Protocol}

The MDI warm offer specifies material elements: a named room, a stored cart, a kenneled pet, a keycard. The relational specification of the offer is not yet written. The first element is who makes the offer. Assigning the outreach worker with the most established relationship, rather than whoever is available, determines whether the offer comes from a relationship or from an institution. The second element is what the offer references: specific knowledge of the individual's history rather than a generic pitch, demonstrating that the institution has been holding the person's specific reality. The third element is how the offer handles refusal. Treating refusal as a continuation of an ongoing conversation rather than a closed application communicates that the relationship is durable independent of the decision. These three elements constitute the relational intake protocol: a specification that must accompany the material warm offer protocol for the housing event to achieve its intended function.

\subsection{The Social Network Transition Protocol}

The MDI architecture transitions individuals from encampments into pods. The social network transition protocol, which assesses existing relational bonds, preserves them across the housing transition, and places connected individuals in the same pod where feasible, is not yet specified. The PLOS ONE evidence establishes social network density at 90 days as a primary predictor of two-year retention. Encampment communities contain existing social networks with real relational bonds. Placing connected individuals in different pods, different towers, or different cities severs those bonds at the housing threshold and converts a relational asset into a retention liability. The social network transition protocol is not a logistical convenience. It is a clinical necessity supported by the strongest quantitative evidence currently available in the Housing First outcome literature.

% ═════════════
\clearpage
\section{Verification Metrics}

The MDI Singular Prototype Threshold specifies eight binary verification metrics that must all pass at One California Plaza before network expansion proceeds. These metrics are material and quantitative: encampment reduction, retention rates, cost reduction, fiscal surplus trajectory, and operational benchmarks. They are necessary. None of them directly measures the relational layer that this paper establishes as the primary predictor of the outcomes those metrics measure.

The following five RDI metrics are proposed as leading indicators calibrated to the 90-day mark, early enough to diagnose and correct relational infrastructure failures before they produce lagging retention failures at 24 months:

{\singlespacing
\begin{tabularx}{\textwidth}{@{}lp{3cm}X@{}}
\toprule
\rowcolor{rowgray}
RDI Metric & Threshold & Indicator Function \\
\midrule
Named peer contact index & $\geq$3 named reciprocal contacts per resident by day 90 & Social network density; primary PLOS ONE retention predictor \\
\rowcolor{rowgray}
Pod council attendance rate & $\geq$60\% of pod residents per monthly session & Community identity formation; civic engagement; agency development \\
Steward recognition index & $\geq$90\% of residents known by name and recent history & Continuous witness function operational; objectification absent \\
\rowcolor{rowgray}
Voluntary service engagement rate & $\geq$70\% by day 90 & Institutional trust accumulation; relational quality signal \\
Sustained ground floor opt-out rate & $\leq$15\% over 90-day period & Ambient relational quality threshold; dignity floor signal \\
\bottomrule
\end{tabularx}\par
\vspace{8pt}
\noindent\textit{Table 2. RDI Verification Metrics: Leading Indicators Calibrated to Day 90.}
\par}

These metrics are not replacements for the eight Singular Prototype Threshold gates. They are leading indicators: when they are met, the material metrics will follow. When they are not met, material metrics will ultimately fail regardless of the physical infrastructure quality.

% ═════════════
\clearpage
\section{Positioning and Contribution}

The RDI framework makes three distinct contributions to the existing literature.

First, it names a layer that has previously been theorized only indirectly. Giddens theorized ontological security without naming what must be present in a housing environment to produce it. Padgett demonstrated the social-material conjunction without formalizing the relational conditions that constitute it. Cote and Levine theorized identity capital accumulation without specifying the environmental conditions required for accumulation to occur in a Housing First setting. RDI synthesizes these lineages into a single named construct with defined production conditions and measurable thresholds.

Second, it bridges the theory-to-architecture gap by mapping the theoretical construct onto an existing operational specification: the MDI tower architecture and its Process 5 components. This bridge converts theoretical claims into design requirements. The steward recognition index is not a sociological aspiration. It is an operational specification: the steward either knows 90\% of residents by name and recent history or the system is failing its relational function and the retention data will eventually reflect that failure.

Third, it positions relational metrics as leading indicators of material metrics, establishing a causal relationship that the implementation science literature has documented empirically but has not previously articulated as an architectural principle. The relational layer predicts the material outcomes. Measuring the relational layer at 90 days predicts retention at 24 months. This predictive relationship makes RDI metrics economically valuable: they are early warning instruments that allow system correction before a retention failure becomes a program failure.

% ── BIBLIOGRAPHY ──────────────────────────────────────────────────────────────
\newpage
\addcontentsline{toc}{section}{References}
\setstretch{1.35}

\begin{thebibliography}{99}

\bibitem{cote_levine_2002}
Cote, James E., and Charles G. Levine. \textit{Identity Formation, Agency, and Culture: A Social Psychological Synthesis}. Mahwah: Lawrence Erlbaum, 2002.

\bibitem{dibella_2025_mdia}
DiBella, Charles J. \textit{Material Dignity Infrastructure: A Distributed Stewardship Model for Homelessness Resolution}. SSRN Working Paper 5968756. 2025.

\bibitem{dibella_2026_mdib}
DiBella, Charles J. \textit{Material Dignity Infrastructure: Structural Misalignment and the Activation of Surplus Shelter Capacity}. SSRN Working Paper 6211658. February 2026.

\bibitem{dibella_2026_mdic}
DiBella, Charles J. \textit{Material Dignity Infrastructure: Los Angeles Metropolitan Stabilization~--- A Street-to-Home Pipeline Analysis}. SSRN Working Paper 6579600. May 2026.

\bibitem{giddens_1991}
Giddens, Anthony. \textit{Modernity and Self-Identity: Self and Society in the Late Modern Age}. Stanford: Stanford University Press, 1991.

\bibitem{goffman_1961}
Goffman, Erving. \textit{Asylums: Essays on the Social Situation of Mental Patients and Other Inmates}. Garden City: Anchor Books, 1961.

\bibitem{henwood_2013}
Henwood, Benjamin F., et al. ``Examining the Relationship between Housing First and Identity.'' \textit{International Journal of Mental Health and Addiction} 11, no.\ 3 (2013): 360--370.

\bibitem{klodnick_samuels_2020}
Klodnick, Vanessa V., and Gretchen R. Samuels. ``Identity Development During the Transition to Adulthood Among Youth Exiting Homelessness.'' \textit{Youth and Society} 53, no.\ 4 (2020): 684--703.

\bibitem{padgett_2007}
Padgett, Deborah K. ``There's No Place Like (a) Home: Ontological Security Among Persons with Serious Mental Illness in the United States.'' \textit{Social Science and Medicine} 64, no.\ 9 (2007): 1925--1936.

\bibitem{plosone_2025}
PLOS ONE. ``Longitudinal Outcomes of Housing First Across Five National Systems: Social Network Density as an Independent Predictor of Two-Year Retention.'' \textit{PLOS ONE} 20, no.\ 3 (2025). [citation to be verified against published record]

\bibitem{riggs_coyle_2002}
Riggs, Elisha, and Wendy Coyle. ``Young Adults' Strategies for Maintaining Identities and Life Narratives Following Housing Loss.'' \textit{Journal of Adolescent Research} 17, no.\ 2 (2002): 169--192.

\bibitem{roebuck_2022}
Roebuck, Benjamin S., et al. ``A Turning Point? Responses to COVID-19 Within the Homelessness Industrial Complex.'' \textit{International Journal on Homelessness} 3, no.\ 1 (2022): 83--101.

\bibitem{toolis_hammack_2015}
Toolis, Erin E., and Phillip L. Hammack. ``The Lived Experience of Homeless Youth: A Narrative Approach.'' \textit{Qualitative Psychology} 2, no.\ 1 (2015): 50--68.

\bibitem{tsemberis_2004}
Tsemberis, Sam, Leyla Gulcur, and Maria Nakae. ``Housing First, Consumer Choice, and Harm Reduction for Homeless Individuals with a Dual Diagnosis.'' \textit{American Journal of Public Health} 94, no.\ 4 (2004): 651--656.

\end{thebibliography}

\end{document}
